%% kapa: exam=1 round=31 year=2020 period=1 form=A subject=civil_law subject_ko=민법 seq=20 number=20
%% source: Q-Net 감정평가사 기출문제 (cst00309) artlSeq=5207219 file=extracted/1008/5207219/A형/A형/2020년도 제31회 감정평가사 1차 시험_1교시_A형.pdf
%% conversion: kordoc 4.24.51 -> tools/gen1.py (auto); review=proofread-sonnet; flags=none
\begin{question}{20}
소멸시효이익의 포기에 관한 설명으로 옳지 않은 것은? (다툼이 있으면 판례에 따름)
\begin{choices}
  \item 시효이익은 미리 포기하지 못한다.
  \item 금전채무에 대한 시효이익의 포기는 채무 전부에 대하여 하여야 한다.
  \item 시효이익을 포기한 때부터 새로이 소멸시효가 진행한다.
  \item 시효이익의 포기는 철회하지 못한다.
  \item 채권의 시효이익을 포기한 경우, 이는 채권자와 채무자의 관계에서만 효력이 생긴다.
\end{choices}
\end{question}
\answer{2}
